% Part: first-order-logic
% Chapter: beyond
% Section: higher-order-logic

\documentclass[../../../include/open-logic-section]{subfiles}

\begin{document}

\olfileid[fr]{fol}{byd}{hol}

\olsection{Logique d'ordre supérieur}

Le passage de la logique du premier ordre à celle du second ordre
nous a permis de parler, dans le langage formel, d'ensembles d'objets
du !!{domain} du premier ordre. Pourquoi s'arrêter là ? La logique du
troisième ordre devrait, par exemple, nous permettre de traiter
d'ensembles d'ensembles d'objets, voire d'ensembles contenant à la
fois des objets et des ensembles d'objets. La logique du quatrième
ordre nous permettrait alors de parler d'ensembles d'objets de cette
sorte. On peut, comme on s'en doute, itérer l'idée à volonté.

En pratique, la logique d'ordre supérieur est souvent formulée en
termes de fonctions plutôt que de relations. (À des identifications
naturelles près, la différence est sans importance.) Étant donné des
« sortes » de base $A$, $B$, $C$,~\dots, que nous appellerons
désormais des « types », nous pouvons en former de nouveaux en posant :
\begin{quote}
Si $\sigma$ et $\tau$ sont des types finis, alors $\sigma \to \tau$
en est un aussi.
\end{quote}
Les types sont des « étiquettes » syntaxiques qui classent les objets
que nous souhaitons inclure dans notre !!{domain}; $\sigma \to \tau$
désigne les objets qui sont des fonctions prenant des objets de
type~$\sigma$ comme arguments et produisant des objets de type~$\tau$.
Nous pourrions, par exemple, disposer d'un type $\Omega$ de valeurs
de vérité, « vrai » et « faux », et d'un type $\Nat$ de nombres
naturels. On peut alors considérer les objets de type $\Nat \to
\Omega$ comme des relations unaires, ou des sous-ensembles de $\Nat$;
ceux de type $\Nat \to \Nat$ comme des fonctions des nombres naturels
vers eux-mêmes; et ceux de type $(\Nat \to \Nat) \to \Nat$ comme des
« fonctionnelles », c'est-à-dire des fonctions de type supérieur qui
prennent des fonctions pour arguments et renvoient des nombres.

Comme pour la logique du second ordre, on peut envisager la logique
d'ordre supérieur comme une logique multi-sortée, avec une sorte pour
chaque type d'objet considéré. Il est toutefois généralement plus
clair de définir directement la syntaxe d'ordre supérieur. Définissons
par exemple, par induction, l'ensemble des types finis :
\begin{enumerate}
\item $\Nat$ est un type fini.
\item Si $\sigma$ et $\tau$ sont des types finis, alors $\sigma
  \to \tau$ en est un aussi.
\item Si $\sigma$ et $\tau$ sont des types finis, alors $\sigma \times
  \tau$ en est un aussi.
\end{enumerate}
Intuitivement, $\Nat$ désigne le type des nombres naturels,
$\sigma \to \tau$ celui des fonctions de $\sigma$ vers $\tau$, et
$\sigma \times \tau$ celui des couples formés d'un objet de $\sigma$
et d'un objet de $\tau$. Définissons ensuite, toujours par induction,
un ensemble de termes :
\begin{enumerate}
\item Pour chaque type $\sigma$, on dispose de variables $x$, $y$,
  $z$, \dots de type $\sigma$.
\item $\Obj 0$ est un terme de type $\Nat$.
\item $\Obj S$ (successeur) est un terme de type $\Nat \to \Nat$.
\item Si $s$ est un terme de type $\sigma$, et $t$ un terme de type
  $\Nat \to (\sigma \to \sigma)$, alors $\Obj{R}_{st}$ est un terme
  de type $\Nat \to \sigma$.
\item Si $s$ est un terme de type $\tau \to \sigma$ et $t$ un terme
  de type~$\tau$, alors $s(t)$ est un terme de type $\sigma$.
\item Si $s$ est un terme de type~$\sigma$ et $x$ une variable de
  type~$\tau$, alors $\lambd[x][s]$ est un terme de type
  $\tau \to \sigma$.
\item Si $s$ est un terme de type~$\sigma$ et $t$ un terme de
  type~$\tau$, alors $\tuple{s, t}$ est un terme de type
  $\sigma \times \tau$.
\item Si $s$ est un terme de type~$\sigma \times \tau$, alors
  $p_1(s)$ est un terme de type~$\sigma$ et $p_2(s)$ un terme de
  type~$\tau$.
\end{enumerate}
Intuitivement, $\Obj{R}_{st}$ désigne la fonction définie par
récurrence au moyen de
\begin{align*}
\Obj{R}_{st}(0) & = s \\
\Obj{R}_{st}(x+1) & = t(x)(\Obj{R}_{st}(x)),
\end{align*}
$\tuple{s, t}$ désigne le couple dont la première composante
est~$s$ et la seconde~$t$, tandis que $p_1(s)$ et~$p_2(s)$ désignent
les première et seconde composantes (les « projections ») de~$s$.\footnote{Précisions éditoriales : « type fini » désigne ici un type engendré par un nombre fini d'étapes syntaxiques, sans affirmer que son domaine est fini. L'étape du récurseur est écrite avec les applications successives requises par le type déclaré de $t$, là où la source écrit une application à deux arguments.}
Enfin, $\lambd[x][s]$ désigne la fonction~$f$ définie par
\[
f(x) = s
\]
pour tout~$x$ de type~$\tau$\footnote{Correction éditoriale : le type de $x$ est $\tau$, comme dans la clause (6), et non $\sigma$, comme l'indique ici la source.}; la clause (6) fournit donc une forme
de compréhension, qui permet de définir des fonctions à l'aide de
termes. Les !!{formula}s se construisent à partir des énoncés
d'!!{identity} $\eq[s][t]$ entre termes de même type, des connecteurs
propositionnels usuels et de quantificateurs portant sur les types
supérieurs. On peut alors prendre pour axiomes les équations
fondamentales qui régissent les termes ainsi définis, accompagnées
des règles logiques usuelles pour les quantificateurs et
l'!!{identity}.

Si l'on ajoute au système des types finis un type $\Omega$ de
valeurs de vérité, il faut également des axiomes régissant son emploi.
Avec un choix judicieux de définitions, on peut même se passer
entièrement des !!{formula}s complexes et les remplacer par des termes
de type $\Omega$ !{} On modifie alors le système de !!{derivation}
en conséquence. On obtient essentiellement la \emph{théorie simple
des types} publiée par Alonzo Church en 1940.\footnote{Correction éditoriale : la source indique les années 1930 ; l'article de Church, \emph{A Formulation of the Simple Theory of Types}, paraît dans le \emph{Journal of Symbolic Logic}, 5(2), en juin 1940, p.~56--68, \href{https://doi.org/10.2307/2266170}{doi:10.2307/2266170}.}

Comme pour la logique du second ordre, plusieurs versions de la
sémantique d'ordre supérieur sont possibles. Dans la version pleine,
les variables de type $\sigma \to \tau$ parcourent l'ensemble de
\emph{toutes} les fonctions des objets de type~$\sigma$ vers ceux de
type~$\tau$. Cette sémantique est trop forte pour
admettre un système de !!{derivation} à la fois complet et effectif.
On peut cependant considérer une sémantique plus faible, où
!!a{structure} se compose d'ensembles d'éléments $T_\tau$ pour
chaque type $\tau$, ainsi que d'opérations appropriées
d'application, de projection, etc. Si l'on précise convenablement
ces notions, on obtient des théorèmes de complétude pour les
systèmes de !!{derivation} décrits plus haut.

La logique d'ordre supérieur est attrayante parce qu'elle fournit un
cadre où une bonne part des mathématiques s'exprime naturellement :
à partir de $\Nat$, on peut définir les nombres réels, les fonctions
continues, et ainsi de suite. Elle présente un intérêt particulier
dans le contexte de la logique intuitionniste, puisque les types y
admettent des interprétations « constructives » claires. On peut
en effet élaborer des versions constructives de la sémantique d'ordre
supérieur, fondées sur la logique intuitionniste plutôt que classique;
elles éclairent ces interprétations et sont, à bien des égards, plus
intéressantes que leurs homologues classiques.

\end{document}
